% =============================================================================
% GRPO-REGISTRY: A Living Catalog of Stack Fields and Variant Deltas
% =============================================================================
% Pillar 1 of the ZVF Program. Companion paper to MIN-REPORT-RL (position).
%
% This is a standalone draft. It is self-contained and should compile with a
% plain `article` class (no external .sty or .bib required: citations are
% rendered as inline TODO keys via a redefined \cite, and the bibliography is a
% manual thebibliography stub). When promoting to a venue template, swap the
% preamble for the venue .sty + \usepackage[numbers]{natbib} and replace TODO
% keys with real \bib entries.
% =============================================================================

\documentclass[11pt]{article}

\usepackage[utf8]{inputenc}
\usepackage[T1]{fontenc}
\usepackage[margin=1in]{geometry}
\usepackage{hyperref}
\usepackage{url}
\usepackage{booktabs}
\usepackage{amsmath}
\usepackage{amssymb}
\usepackage{xcolor}
\usepackage{enumitem}
\usepackage{microtype}
\usepackage{graphicx}
\usepackage{longtable}
\usepackage{listings}
\usepackage{algorithm}
\usepackage{algpseudocode}

% ---------------------------------------------------------------------------
% TODO-citation shim. We do NOT invent bib entries with fake authors/years.
% Every \cite below uses a clearly-TODO key and renders in red so a human can
% find and replace it with a real reference + .bib entry before submission.
% ---------------------------------------------------------------------------
\renewcommand{\cite}[1]{\textcolor{red}{[\textsc{cite:}\,#1]}}
\newcommand{\todo}[1]{\textcolor{red}{\textbf{[TODO:\ #1]}}}
\newcommand{\zvf}{\textsc{Zvf}}
\newcommand{\gu}{\textsc{Gu}}
\newcommand{\minreport}{\textsc{Min-Report-RL}}
\newcommand{\grpo}{\textsc{GRPO}}
\newcommand{\grpoReg}{\textsc{GRPO-Registry}}

% JSON listing style
\lstdefinelanguage{json}{
  basicstyle=\ttfamily\small,
  numbers=left,
  numberstyle=\tiny\color{gray},
  stepnumber=1,
  numbersep=8pt,
  showstringspaces=false,
  breaklines=true,
  frame=single,
  string=[s]{"}{"},
  comment=[l]{//},
  morecomment=[s]{/*}{*/},
  keywordstyle=\color{blue}\bfseries,
  stringstyle=\color{teal},
  literate={:}{{\color{orange}:}}1
           {,}{{\color{orange},}}1
           {\{}{{\color{purple}\{}}1
           {\}}{{\color{purple}\}}}1
           {[}{{\color{purple}[}}1
           {]}{{\color{purple}]}}1,
}

\title{\grpoReg{}: A Living Catalog of Stack Fields and Variant Deltas\\
\large Toward Machine-Readable Reporting for the \grpo{} Family}

\author{Arvind C R \and the ZVF Program\thanks{Companion to the ZVF Program
  \minreport{} position paper. Author list and affiliations \todo{finalize}.}}

\date{\today}

\begin{document}
\maketitle

% =============================================================================
\begin{abstract}
% =============================================================================
The \grpo{} family of reinforcement-learning post-training methods is now
large enough that no single paper can keep its variants straight. Each new
variant---DAPO, GSPO, Dr.\grpo{}, MAD-\grpo{}, AERO, CPPO, NGRPO,
Scaf-\grpo{}, and others---is reported as a small algorithmic delta, but the
delta is almost always confounded with changes in loss form, sampler, reference
policy / KL handling, group-size schedule, and precision. We introduce
\grpoReg{}, a living catalog that (i) fixes a common seven-item
minimum-reportable-stack schema drawn from the \minreport{} standard, (ii)
defines a machine-readable ``variant-delta'' schema that records exactly what
changed relative to a declared baseline, and (iii) seeds the catalog with the
best currently available descriptions of the most discussed variants. The goal
is not to arbitrate priority but to give reviewers and re-implementers a shared
worksheet: for any comparison, read the two stack records, read the two variant
records, and see whether the claimed gain is still attributable once the stack
is held fixed. We release the schema as JSON and propose a lightweight
update protocol so the catalog can stay current without a new paper for every
new variant.
\end{abstract}

% =============================================================================
\section{Introduction: Why a Registry?}
\label{sec:intro}
% =============================================================================

A reviewer reading six recent \grpo{} papers can reasonably conclude that each
paper compares against ``GRPO.'' But the word ``GRPO'' is attached to different
loss forms (ratio + clip, ratio + no clip, completion-only mask,
whole-sequence mask), different reference-policy and KL treatments (frozen
reference with KL penalty in the loss, frozen reference with KL in the reward,
no reference at all), different samplers and precisions (vLLM bf16, vLLM fp32,
managed API), and different group-size schedules (fixed $G=8$, adaptive,
dynamic sampling) \cite{zvfaudit2026, dapo2025, gspo2025, drgrpo2025,
madgrpo2025}. When the baseline is itself a moving target, a reported delta of
$+2$ points is uninterpretable.

The position paper \minreport{} \cite{minreportrl2026} argues that the first
fix is a short, mandatory list of stack fields every paper must report. That
fixes the ``what is the treatment?'' problem for a single paper. It does not,
by itself, fix the cross-paper comparison problem: a reader still has to open
every paper, extract the seven fields by hand, and decide which variant is
which. \grpoReg{} is that extraction, maintained as a shared artifact.

\paragraph{What the registry is.} \grpoReg{} has three parts:
\begin{enumerate}[leftmargin=1.6em, itemsep=2pt]
  \item \textbf{A stack schema} (\S\ref{sec:stack-schema}), the seven
    \minreport{} fields, expressed as a JSON object any trainer can emit.
  \item \textbf{A variant-delta schema} (\S\ref{sec:delta-schema}) that
    records, relative to a named baseline, which knobs the variant touches.
  \item \textbf{A catalog of variants} (\S\ref{sec:catalog}) populated with
    the best available descriptions of the major \grpo{} variants; cells whose
    exact setting is unknown are left as \todo{} placeholders.
\end{enumerate}

\paragraph{What the registry is not.} It is not a leaderboard. It does not
score variants. It is a structured worksheet for de-confounding comparisons:
aid the reader in asking ``did method $A$ beat method $B$, or did stack $S_A$
beat stack $S_B$?'' \cite{zvfaudit2026}.

% =============================================================================
\section{The Stack Schema}
\label{sec:stack-schema}
% =============================================================================

The stack schema is the seven-item \minreport{} checklist \cite{minreportrl2026},
reformatted as named JSON fields. Table~\ref{tab:stack-schema} gives the
human-readable mapping. The key design choice is that each field is included
because it is a documented lever that can flip a comparison; the justification
is in the position paper and summarized here.

\begin{table}[t]
\centering
\caption{The \grpoReg{} stack schema. Each field is one item of the
  \minreport{} minimum-reportable-stack standard.}
\label{tab:stack-schema}
\small
\begin{tabular}{@{}lp{0.55\linewidth}l@{}}
\toprule
\textbf{JSON key} & \textbf{What to report} & \textbf{Flip mechanism}\\
\midrule
\texttt{loss\_form} & Ratio? Clip? Asymmetric clip? Token mask? Advantage norm.?
  & Reassigns gradient; DAPO, GSPO, Dr.\grpo{} deltas live here.\\
\texttt{reference\_kl} & Frozen ref? KL in loss / reward / absent? Coefficient &
  Changes objective and exploration; no-ref vs.\ KL are different methods.\\
\texttt{sampler\_backend} & Engine, decoding params, precision, tokenizer match &
  Backend/precision drove a $\sim$17$\times$ matched-config gap
  \cite{zvfaudit2026}.\\
\texttt{zvf\_gu} & Per-step \zvf{} and \gu{} trajectory &
  Reveals whether the run had usable group-relative signal.\\
\texttt{group\_schedule} & Fixed or adaptive $G$; if adaptive, rule &
  $G$ controls mixed-group probability; adaptive-$G$ confounds compute.\\
\texttt{held\_out} & Disjoint eval slice, $n$, CI, reported separately &
  Training reward is dynamics evidence, not capability evidence.\\
\texttt{decontamination} & Overlap probe + parser behavior on format-only inputs &
  Verifiable rewards still admit parser/format/overlap hacking.\\
\bottomrule
\end{tabular}
\end{table}

The schema is intentionally small. A trainer that already logs its config and
per-step group rewards can emit it with a few lines of code. The companion
\minreport{} paper proposes default emitters for TRL, verl, and OpenRLHF
\cite{minreportrl2026}.

% =============================================================================
\section{The Variant-Delta Schema}
\label{sec:delta-schema}
% =============================================================================

The variant-delta schema records what changed relative to a declared baseline.
A delta record is 	extbf{not} a full stack record; it is a patch. For example,
DAPO's delta record says ``relative to a vanilla \grpo{} baseline, change the
clip bound to be asymmetric and add dynamic sampling,'' while the full stack
record still has to report the sampler, the reference/KL handling, the
precision, and so on. This separation keeps the catalog compact and forces
authors to state both the variant and the stack it ran on.

Table~\ref{tab:delta-schema} lists the delta fields. They are chosen to cover
the dimensions on which the current \grpo{} variants actually differ.

\begin{table}[t]
\centering
\caption{The \grpoReg{} variant-delta schema. Each field records a change
  relative to the variant's declared baseline. ``\texttt{null}'' means the
  variant does not change that field.}
\label{tab:delta-schema}
\small
\begin{tabular}{@{}lp{0.65\linewidth}@{}}
\toprule
\textbf{JSON key} & \textbf{Meaning}\\
\midrule
\texttt{variant\_name} & Canonical short name (e.g., ``DAPO'').\\
\texttt{baseline} & The baseline the variant claims to improve (usually
  ``GRPO'').\\
\texttt{loss\_change} & Change to the objective (ratio, clip, token mask,
  normalization).\\
\texttt{clip\_rule} & New clip bound or asymmetric rule, if any.\\
\texttt{sampling\_change} & Change to rollout sampling or group construction
  (dynamic, filter, \ldots).\\
\texttt{group\_size\_rule} & Fixed / adaptive; if adaptive, the rule.\\
\texttt{advantage\_norm} & How advantages are normalized (per-group,
  per-batch, running).\\
\texttt{reference\_kl\_change} & Any change to reference policy or KL handling.\\
\texttt{extra\_regularizer} & Additional penalty or bonus (length, diversity,
  entropy, \ldots).\\
\texttt{rollout\_change} & Change to generation parameters or decoding.\\
\texttt{other\_notes} & Anything else that defines the variant.\\
\bottomrule
\end{tabular}
\end{table}

A variant record is therefore a pair:
\[
\text{variant record} = \langle \text{stack schema}, \text{delta schema} \rangle .
\]
A fair head-to-head between two variants is a comparison of two such pairs in
which the stack schemas are identical except where the delta explicitly changes
them.

% =============================================================================
\section{Catalog of \grpo{} Variants}
\label{sec:catalog}
% =============================================================================

Table~\ref{tab:catalog} is the current seed catalog. We include the best
available description of each variant's defining delta, but we do 	extbf{not}
fabricate stack details. Unknown stack settings are marked \todo{} so that a
reader can see exactly what is still unspecified. The catalog is versioned;
as papers release more detail, \todo{} cells are replaced with concrete values.

\begin{longtable}{@{}p{2.1cm}p{2.4cm}p{3.0cm}p{2.0cm}p{3.3cm}@{}}
\caption{Seed catalog of \grpo{}-family variants. Each row describes the
  variant's delta relative to its baseline. Empty or \todo{} cells indicate
  settings not yet reported in the source literature.}\\
\toprule
\textbf{Variant} & \textbf{Baseline} & \textbf{Loss / clip change} &
\textbf{Sampling / $G$ change} & \textbf{Extra regularizer / notes}\\
\midrule
\endfirsthead
\multicolumn{5}{c}{\tablename\ \thetable{} -- continued}\\
\toprule
\textbf{Variant} & \textbf{Baseline} & \textbf{Loss / clip change} &
\textbf{Sampling / $G$ change} & \textbf{Extra regularizer / notes}\\
\midrule
\endhead
\midrule
\multicolumn{5}{r}{\textit{Continued on next page}}\\
\endfoot
\bottomrule
\endlastfoot

\grpo{} (baseline) & --- & Per-token GRPO ratio with clip; completion-only
  mask \todo{verify} & Fixed group $G$ & KL term placement varies by
  implementation \cite{zvfaudit2026}.\\
\addlinespace

DAPO & \grpo{} & Asymmetric clip-higher lower bound & Dynamic sampling:
  group constructed from filtered / top-$k$ completions & Adds per-token
  loss-clip asymmetry \cite{dapo2025}.\\
\addlinespace

GSPO & \grpo{} & Sequence-level importance-sampling ratio rather than
  per-token ratio; clip rule \todo{} & Group-size rule \todo{} & Changes
  importance-sampling granularity \cite{gspo2025}.\\
\addlinespace

Dr.\grpo{} & \grpo{} & Length-normalized advantage and reward terms; fixes
  length bias in the per-token loss & Fixed or adaptive $G$ \todo{} & Removes
  length bias in normalization \cite{drgrpo2025}.\\
\addlinespace

MAD-\grpo{} & \grpo{} & Adds multi-agent / diversity term into the group
  advantage or reward & Group construction may be diversity-aware & Diversity
  bonus across agents \cite{madgrpo2025}.\\
\addlinespace

AERO & \grpo{} & Usually unchanged \todo{} & Adaptive group-size / rollout
  length based on a rolling signal estimate & Spend-more-compute-on-hard-prompts
  rule \cite{aero2024}.\\
\addlinespace

CPPO & PPO / \grpo{} & Clip-pruned objective: tokens with large ratios are
  dropped instead of clipped & Fixed or adaptive \todo{} & Pruning changes the
  effective sample \cite{cppo2024}.\\
\addlinespace

NGRPO & \grpo{} & Normalized advantage / reward transform \todo{specify} &
  Fixed or adaptive \todo{} & Normalization is the entire delta
  \cite{ngrpo2025}.\\
\addlinespace

Scaf-\grpo{} & \grpo{} & Incorporates scaffolded exploration signal into the
  advantage \todo{specify} & Fixed or adaptive \todo{} & Scaffolded hints shape
  the rollout distribution \cite{scafgrpo2025}.\\
\addlinespace

GRESO & \grpo{} & Unchanged \todo{verify} & \emph{Pre-rollout} filtering:
  prompts predicted to yield zero-gradient (single-outcome) groups are skipped
  before any rollout is spent & Acts on the dead-group phenomenon \emph{before}
  sampling; the prediction rule is itself a stack component
  \cite{greso2025,agenticrl2026}.\\
\addlinespace

EDGE-\grpo{} & \grpo{} & Entropy-driven advantage reshaping + guided error
  correction & Fixed $G$ \todo{verify} & Explicitly framed as mitigating
  \emph{advantage collapse} --- the same phenomenon \zvf{} measures
  \cite{edgegrpo2025,agenticrl2026}.\\
\addlinespace

DARS & \grpo{} & Unchanged \todo{verify} & Multi-stage rollout sampling that
  reallocates compute from medium-difficulty to the hardest prompts &
  Difficulty-aware rollout budget; confounds compute allocation with the
  algorithmic claim unless reported \cite{dars2025,agenticrl2026}.\\
\addlinespace

TreePo & \grpo{} & Unchanged \todo{verify} & Self-guided policy rollout over a
  tree structure to reduce rollout compute & Compute-reduction delta; group
  construction differs from i.i.d.\ sampling \cite{treepo2025,agenticrl2026}.\\
\addlinespace

Posterior-\grpo{} & \grpo{} & Rewards only successful \emph{processes}
  (intermediate steps), not outcomes alone & Fixed \todo{verify} & Changes the
  reward density and hence steady-state \zvf{}
  \cite{posteriorgrpo2025,agenticrl2026}.\\
\addlinespace

Skywork-R1V2 & \grpo{} & Hybrid reward signal \todo{specify} & Selective
  sample buffer replays informative (mixed-reward) samples & Buffer selection
  is a direct intervention on group reward diversity
  \cite{skyworkr1v2-2025,agenticrl2026}.\\
\addlinespace

GMPO & \grpo{} & Geometric mean of token-level rewards in place of arithmetic
  aggregation & Fixed \todo{verify} & Aggregation change alters which groups
  register as zero-variance \cite{gmpo2025,agenticrl2026}.\\
\addlinespace

ProRL & \grpo{} & Unchanged \todo{verify} & Fixed \todo{verify} &
  Reference-policy \emph{reset} during prolonged training; changes KL anchor
  identity mid-run --- a stack field the label omits
  \cite{prorl2025,agenticrl2026}.\\

\end{longtable}
\label{tab:catalog}

\paragraph{Reading the table.} A cell marked \todo{} is not a weakness of the
registry; it is a weakness of the underlying paper's reporting. The registry
surfaces it so that a re-implementer knows what still has to be reverse
engineered. A fair comparison between, say, DAPO and Dr.\grpo{} requires not
only the delta rows above but also the full stack records of the two original
runs; only then can one tell whether the reported gap is a stack effect.

\paragraph{The dead-group lever is now the most-attacked knob in the family.}
The TMLR survey of agentic RL \cite{agenticrl2026} canonizes, in a single
comparison table, more than twenty \grpo{}-family variants --- and a striking
fraction of them (DAPO's dynamic sampling, GRESO's pre-rollout filtering,
EDGE-\grpo{}'s \emph{advantage collapse} mitigation, DARS's difficulty-aware
reallocation, Skywork-R1V2's selective buffer, TreePo's rollout pruning)
intervene, each with a different heuristic, on precisely the phenomenon
\zvf{} measures: groups whose identical rewards contribute zero gradient.
We take the survey's adoption of the term \emph{advantage collapse} as
confirmation that the field has converged on the phenomenon; what it has not
converged on is a calibrated way to \emph{measure} it, which is exactly the
telemetry this registry and the \minreport{} standard demand. None of the
variant papers above reports a per-step dead-group trajectory, so their
relative gains on this axis are currently unattributable.

% =============================================================================
\section{How to Use the Registry to Flag Confounded Comparisons}
\label{sec:usage}
% =============================================================================

The registry is a worksheet, not an oracle. The recommended workflow for a
reviewer or re-implementer is:

\begin{enumerate}[leftmargin=1.6em, itemsep=2pt]
  \item \textbf{Extract the stack record} for each paper. If a paper does not
    report one of the seven fields, mark it \todo{}.
  \item \textbf{Extract the variant-delta record} for the method claimed. If
    the paper does not explicitly state what changed relative to its baseline,
    the delta is itself \todo{}.
  \item \textbf{Align baselines.} The comparison is only meaningful if both
    variants declare the same baseline (usually vanilla \grpo{}).
  \item \textbf{Check for stack confounds.} Any stack field that differs
    between the two papers and is not part of either variant's declared delta
    is a confound.
  \item \textbf{Check for delta overlap.} If variant $A$ changes the clip rule
    and variant $B$ also changes the clip rule, a head-to-head is comparing two
    clip rules, not ``method $A$ vs.\ method $B$.''
  \item \textbf{Look at \zvf{} / \gu{}.} If one paper reports high \zvf{}
    (low usable signal) and the other reports low \zvf{}, the comparison is
    confounded by signal availability, not algorithmic quality.
\end{enumerate}

\subsection{Mechanical checks with \texttt{grpo-stackdiff}}

The registry's records are designed to be consumed by the companion
\texttt{grpo-stackdiff} CLI \cite{minreportrl2026}. Given two
\grpoReg{} JSON records, \texttt{grpo-stackdiff} classifies each field
difference by lever (L1--L7), diff kind, role (treatment, nuisance, coverage
gap, invalidator), and flip risk (R0--RU). A comparison whose maximum nuisance
risk is R4 (``flip-capable'') or R5 (``invalidating'') is mechanically
flagged as \texttt{STACK\_CONFOUNDED} or \texttt{INVALID\_COMPARISON}. The
tool therefore turns the six-step manual checklist above into a single
command, and it can be run in CI to block confounded benchmark claims before
they are merged.

A confounded comparison is not wrong; it is underspecified. The registry makes
the underspecification explicit so that the community can decide whether the
available evidence is sufficient.

% =============================================================================
\section{Tooling Ecosystem: From Records to Executable Checks}
\label{sec:tools}
% =============================================================================

The registry is the static worksheet. Three companion tools turn it into an
enforceable workflow.

\paragraph{\texttt{trl-min-report-rl}.} A small TRL extension that emits the
run-start manifest and per-step \zvf{}/\gu{} telemetry automatically. Its
package contains a \texttt{MinReportGRPOTrainer} (or an
\texttt{enable\_min\_report\_rl} adapter for existing trainers), a
reward-tensor hook, and stable JSON schemas for run-start and step events
\cite{minreportrl2026}. The emitted manifest uses exactly the seven-item schema
of \S\ref{sec:stack}, so a training run can populate a registry record with no
manual transcription.

\paragraph{\texttt{grpo-stackdiff}.} The companion CLI takes two registry
records (or two run-start manifests) and classifies every field difference by
lever, diff kind, role, and flip-risk class. It can be run in CI to block a
confounded benchmark claim before it is merged
(\S\ref{sec:mechanical}).

\paragraph{\texttt{MIN-REPORT-RL Auditor}.} For already-published work, the
Auditor scans a paper PDF, code release, and log artifacts; extracts evidence
for the seven \minreport{} items; assigns a 0--10 score per item and a
normalized 0--100 badge (Complete / Usable / Partial / Not Auditable /
Non-compliant); and produces a reviewer-facing summary with quoted evidence,
missing fields, and contradiction warnings \cite{zvfaudit2026}. It is the
post-hoc enforcement layer: it tells authors exactly which fields to add before
the paper is reproducibility-auditable.

\paragraph{\texttt{MRRL-AuditRunner}.} The \minreport{} standard and the
registry together enable a pre-registered controlled audit. The AuditRunner
freezes a shared stack, expands each variant as a declared hook, runs the arms
on paired seeds, and applies survival verdicts from seed-level held-out deltas
(\cite{zvfaudit2026}; full protocol in the audit paper). Its outputs
(Tables~A--H of the runner specification) include the pre-registration bundle,
stack-lock table, arm registry, run ledger, survival table, telemetry summary,
decontamination probes, and artifact index.

Together these tools close the loop: trainers emit manifests, the registry
collects them, \texttt{grpo-stackdiff} checks pairwise comparability, the
Auditor scores legacy papers, and the AuditRunner produces controlled survival
evidence for new claims.

% =============================================================================
\section{Machine-Readable Schema}
\label{sec:json}
% =============================================================================

Listing~\ref{lst:schema} gives the JSON schema for a complete \grpoReg{}
record. A record contains one stack object, one baseline object, and one delta
object. The schema is intentionally flat: it should be easy to generate from a
trainer's config and to diff with standard tools.

\begin{lstlisting}[language=json, caption={JSON schema for a \grpoReg{}
  record. Required fields are marked \texttt{required}; all others may be
  \texttt{null} when not applicable.}, label={lst:schema}]
{
  "version": "grpo-registry-v0.1",
  "record_type": "variant_stack_record",
  "variant": {
    "name": "DAPO",
    "paper_doi_or_url": "TODO",
    "baseline": "GRPO"
  },
  "stack": {
    "loss_form": {
      "ratio": true,
      "clip_bounds": [0.2, 0.2],
      "clip_asymmetric": true,
      "token_mask": "completion-only",
      "advantage_norm": "per-group"
    },
    "reference_kl": {
      "frozen_reference": true,
      "kl_placement": "loss",
      "kl_coefficient": 0.01,
      "kl_schedule": "constant"
    },
    "sampler_backend": {
      "engine": "vLLM",
      "temperature": 1.0,
      "top_p": 1.0,
      "max_tokens": 4096,
      "logit_precision": "bf16",
      "tokenizer_matches_trainer": true
    },
    "zvf_gu": {
      "per_step_zvf": [0.95, 0.92, ...],
      "per_step_gu": [0.05, 0.08, ...]
    },
    "group_schedule": {
      "type": "fixed",
      "fixed_size": 8
    },
    "held_out": {
      "split": "gsm8k-test",
      "n": 1319,
      "metric": "pass@1",
      "value": 0.55,
      "ci_95": [0.52, 0.58]
    },
    "decontamination": {
      "overlap_check": "n-gram-8",
      "overlap_rate": 0.02,
      "parser_probe": "format-only",
      "parser_probe_pass": false
    }
  },
  "delta": {
    "loss_change": "asymmetric clip-higher",
    "clip_rule": "lower=0.2, upper=0.28",
    "sampling_change": "dynamic top-k group construction",
    "group_size_rule": null,
    "advantage_norm": null,
    "reference_kl_change": null,
    "extra_regularizer": null,
    "rollout_change": null,
    "other_notes": "DAPO paper introduces clip-higher + dynamic sampling."
  }
}
\end{lstlisting}

The schema is available at the repository linked in the author notes
\todo{add repository URL}. We welcome pull requests that fill \todo{} cells or
add new variants, provided each entry cites the source paper and does not
invent unreported stack details.

The run-start manifest emitted by \texttt{trl-min-report-rl}
(\S\ref{sec:tools}) is a strict superset of this record: it adds run identity,
software versions, hardware, integrity hashes, and a per-step telemetry stream.
A registry record can therefore be generated automatically from the first line
of \texttt{min\_report\_rl.jsonl}.

% =============================================================================
\section{Living-Update Protocol}
\label{sec:update}
% =============================================================================

A registry that requires a new paper for every new variant will quickly become
stale. We propose a lightweight update protocol:

\begin{enumerate}[leftmargin=1.6em, itemsep=2pt]
  \item \textbf{Versioned releases.} The catalog is released as a versioned
    JSON file (e.g., \texttt{grpo-registry-v0.1.json}). Each release has a
    changelog.
  \item \textbf{Source requirement.} Every entry must cite a specific paper,
    preprint, or code release. No entry may be added purely from a blog post or
    social-media summary.
  \item \textbf{Uncertainty marking.} Unknown fields are recorded as
    \texttt{null} with a \texttt{source\_note} explaining what was not
    reported. This preserves the honesty of the catalog.
  \item \textbf{Community curation.} Changes are proposed via pull request and
    reviewed by at least one maintainer for source fidelity.
\end{enumerate}

The protocol mirrors existing living artifacts such as the Papers With Code
leaderboards, but it is organized around stack fields rather than final scores.

% =============================================================================
\section{Conclusion}
\label{sec:conclusion}
% =============================================================================

The \grpo{} family has outgrown its own label. The registry proposed here is a
small, practical step toward a literature in which the treatment---both the
stack and the declared algorithmic delta---is reported clearly enough to
compare across papers. It does not replace the \minreport{} standard; it
depends on it. It also does not replace controlled re-implementations; it makes
the need for them visible earlier. With the companion tooling of
\S\ref{sec:tools}---the \texttt{trl-min-report-rl} emitter,
\texttt{grpo-stackdiff}, the Auditor, and the AuditRunner---the registry moves
from a static worksheet to an enforceable workflow. The immediate next step is
to populate the catalog with the full stack records of the major variants and to
ship default emitters so that future papers produce a \grpoReg{} record with no
extra writing.

% =============================================================================
% Bibliography. MANUAL STUB ONLY. We do NOT invent authors/years. Every entry
% is a TODO key to be replaced by a real citation before submission.
% =============================================================================
\begin{thebibliography}{9}

\bibitem{zvfaudit2026}
\todo{ZVF Program v1 audit paper. Cite the arXiv/venue version once assigned.}

\bibitem{minreportrl2026}
\todo{\minreport{} position paper. Cite the arXiv/venue version once assigned.}

\bibitem{dapo2025}
Yu, Q., et al. \emph{DAPO: An Open-Source LLM Reinforcement Learning System
at Scale}. arXiv:2503.14476, 2025. (Clip-higher + dynamic sampling.)
\todo{verify final venue}

\bibitem{gspo2025}
Zheng, C., et al. (Qwen Team). \emph{Group Sequence Policy Optimization}.
arXiv:2507.18071, 2025. (Sequence-level importance sampling.)
\todo{verify final venue}

\bibitem{drgrpo2025}
Liu, Z., et al. \emph{Understanding R1-Zero-Like Training: A Critical
Perspective}. arXiv:2503.20783, 2025. (Dr.\grpo{}: length/normalization-bias
fix.) \todo{verify final venue}

\bibitem{madgrpo2025}
\todo{MAD-\grpo{} paper (multi-agent/diversity). Fill authors/venue/year.}

\bibitem{aero2024}
\todo{AERO (adaptive rollout sizing). Fill authors/venue/year.}

\bibitem{cppo2024}
\todo{CPPO (clip-pruned PPO). Fill authors/venue/year.}

\bibitem{ngrpo2025}
\todo{NGRPO (normalized GRPO). Fill authors/venue/year.}

\bibitem{scafgrpo2025}
\todo{Scaf-\grpo{} (scaffolded exploration). Fill authors/venue/year.}

\bibitem{agenticrl2026}
Zhang, G., Geng, H., Yu, X., Yin, Z., et al. \emph{The Landscape of Agentic
Reinforcement Learning for LLMs: A Survey}. Transactions on Machine Learning
Research, 01/2026. arXiv:2509.02547. (Table~2 is the canonizing catalog of
\grpo{}-family variants; source of the variant attributions below.)

\bibitem{greso2025}
Zheng, H., et al. \emph{Act Only When It Pays: Efficient Reinforcement
Learning for LLM Reasoning via Selective Rollout} (GRESO).
arXiv:2506.02177, 2025. \todo{verify final venue}

\bibitem{edgegrpo2025}
Zhang, X., et al. \emph{EDGE-GRPO: Entropy-Driven GRPO with Guided Error
Correction for Advantage Diversity}. arXiv:2507.21848, 2025.
\todo{verify final venue}

\bibitem{dars2025}
Yang, Z., et al. \emph{Depth-Breadth Synergy in RLVR: Unlocking LLM
Reasoning Gains with Adaptive Exploration} (introduces DARS,
difficulty-adaptive rollout sampling). arXiv:2508.13755, 2025.
\todo{verify final venue}

\bibitem{treepo2025}
Li, Y., et al. \emph{TreePO: Bridging the Gap of Policy Optimization and
Efficacy and Inference Efficiency with Heuristic Tree-based Modeling}.
arXiv:2508.17445, 2025. \todo{verify final venue}

\bibitem{posteriorgrpo2025}
Fan, L., Zhang, Y., Chen, M., and Liu, Z. \emph{Posterior-GRPO: Rewarding
Reasoning Processes in Code Generation}. arXiv:2508.05170, 2025.
\todo{verify final venue}

\bibitem{skyworkr1v2-2025}
Wang, P., et al. \emph{Skywork R1V2: Multimodal Hybrid Reinforcement
Learning for Reasoning}. arXiv:2504.16656, 2025. \todo{verify final venue}

\bibitem{gmpo2025}
Zhao, Y., et al. \emph{Geometric-Mean Policy Optimization}.
arXiv:2507.20673, 2025. \todo{verify final venue}

\bibitem{prorl2025}
Liu, M., et al. \emph{ProRL: Prolonged Reinforcement Learning Expands
Reasoning Boundaries in Large Language Models}. arXiv:2505.24864, 2025.
\todo{verify final venue}

\end{thebibliography}

\end{document}
